\section{Separation and the geometry of the accumulation families}
\label{sec:accumulation-geometry}

We first establish the single nontrivial identification for each of the
terminal-state, ordered-path and compact-image parameterizations, then
use continuity and compactness to identify their images with the
quotient circle.

\begin{lemma}[Separation of terminal states and path images]
\label{lem:terminal-separation}
Fix $\lambda,\nu$ and $0<s\le s_0$ as in
Theorem~\ref{thm:whole}, and write $d_*=F_\lambda^{-1}(s)$.
The terminal-state map
\[
 E_s:\mathcal A_s\longrightarrow\mathbb R^5,
 \qquad E_s(d,\alpha)=\cC_{\lambda,d}(\alpha),
\]
is one-to-one except for
\[
 E_s(d_*,3)=E_s(s,0).
\]
The ordered-path and compact-image maps in
\eqref{eq:ps5-whole}--\eqref{eq:ps5-whole-image}
have exactly the same single nontrivial identification.
\end{lemma}

\begin{proof}
The contact injectivity established in
\eqref{gr:contact-injectivity} applies for every $0<s\le s_0$.
It separates admission, the interior contact points and release.

For $d_*<d<s$, the stopped endpoint lies strictly inside its first
free flight and has $y>0$. Contact endpoints have $y=0$.
Two free endpoints cannot coincide. Indeed, suppose
\[
 \Phi_{\rm f}(t_1;\mathcal R_{d_1})
   =\Phi_{\rm f}(t_2;\mathcal R_{d_2}),
 \qquad 0<t_i<T_{d_i},
\]
and take $t_1\le t_2$. Uniqueness for the autonomous free mechanical
flow gives
$\mathcal R_{d_1}=\Phi_{\rm f}(t_2-t_1;\mathcal R_{d_2})$.
If $t_1<t_2$, the right side has $y>0$, whereas the left side has
$y=0$. Hence $t_1=t_2$ and the release states agree. Their $p$ values
give $d_1=d_2$, and strict decrease of $p$ on the free flight gives
the same stopping point. At the free endpoints, graph invariance
gives $\cC_{\lambda,d_*}(3)=\cC_{\lambda,s}(0)$, while
$\cC_{\lambda,s}(2)$ is already the common point of the graph and
vertical branches of $\mathcal A_s$. Injectivity of the contact curve
now proves the assertion for $E_s$.

The two descriptions $(d_*,3)$ and $(s,0)$ give the same ordered
path, since the final return at $d_*$ is the last full return in
$L_s$. Conversely, the ordered distance bounds the distance
between terminal states. The terminal-state assertion therefore gives
the assertion for ordered paths.

For images, endpoint separation alone would not suffice. Put
$\Sigma_s=\{\cX:p(\cX)=-as\}$. The incoming arc has $p\ge0$.
When $d<s$, all completed earlier returns and the final contact have
$p>-as$, and the last free prefix has strictly decreasing $p$.
Consequently
\[
 W_{\lambda,\nu,s}(d,j_s(d))\cap\Sigma_s
          =\{E_s(d,j_s(d))\},\qquad d_*\le d<s.
\]
For a contact prefix, the only earlier point on $\Sigma_s$ is its
admission point: the last full return in $L_s$ ends there, and all
other earlier points have $p>-as$. Hence
\[
 W_{\lambda,\nu,s}(s,\alpha)\cap\Sigma_s
       =\cC_{\lambda,s}([0,\alpha]),\qquad0\le\alpha\le2.
\]
These sets are strictly nested as $\alpha$ increases, by contact
injectivity. A contact prefix of positive length cannot have a
singleton intersection with $\Sigma_s$. The case $\alpha=0$ has
the same singleton as $(d_*,3)$; every interior free endpoint has
$y>0$. Finally two free singletons coincide only in the already
identified terminal-state cases. This proves separation of the
compact images as well.
\end{proof}

\begin{corollary}[The accumulation families are topological circles]
\label{cor:accumulation-circle}
For fixed $\lambda,\nu,s$ as in Lemma~\ref{lem:terminal-separation},
the terminal-state set, the ordered-path accumulation family with
the distance \eqref{eq:ps5-frechet}, and the compact-image accumulation
family with its Hausdorff distance are each homeomorphic to $S^1$.
More precisely, their parameter space is
\[
 \mathcal A_s/\bigl((d_*,3)\sim(s,0)\bigr).
\]
These statements concern the full physical state
$(\phi,y,\omega,v,z)$.
\end{corollary}

\begin{proof}
The graph branch of $\mathcal A_s$ is an arc from $(d_*,3)$ to
$(s,2)$, and its vertical branch is an arc from $(s,2)$ to $(s,0)$.
Their only common point is $(s,2)$, so $\mathcal A_s$ is homeomorphic to a closed
interval with endpoints $(d_*,3)$ and $(s,0)$. Identifying these
endpoints gives $S^1$.

The terminal-state map is continuous by the construction of the
reference curves. The ordered-path map is continuous by the finite-final-part and
small-tail argument following \eqref{path:whole-error}; image
continuity follows from $\dH\le\dF$. Lemma~\ref{lem:terminal-separation} says that all three
maps identify exactly the two interval endpoints. Each therefore
induces a continuous bijection from the compact quotient circle to
its image in a metric space, and is a homeomorphism.
\end{proof}
